\section{Endpoint cutoffs and precise scope}\label{sec:scope}
The main theorem needs only initial compact tests and exact conservation.
The following separate endpoint gate records what Gaussian birth energy
also proves, and the order of limits it permits.

\begin{lemma}[Pasted polynomial endpoint covariance gate]
\label{lem:endpoint}
Let $q$ be any fixed polynomial, $0\le\chi\le1$ a smooth cutoff equal
to one on the unit phase ball and zero outside the radius-two ball,
$\chi_R(z)=\chi(z/R)$, and $q^{>R}=q(1-\chi_R)$. For every fixed
$M>0$, uniformly for sufficiently large $R$ and sufficiently small $\eps$,
\begin{equation}\label{eq:endpoint-tail}
 \sup_{t\le T}\E|\zeta_t^{\eps,p}(q^{>R})|^2
        \le C_qe^{-c_qR^2}+C_{q,M}\eps^M.
\end{equation}
The statement holds for fixed finite endpoint record lists and gives
covariance cutoff approximation at two times. In particular its cutoff
limit takes $\eps\to0$ first, then $R\to\infty$.
\end{lemma}
\begin{proof}
The actual pasted speed bound \eqref{eq:record-clip-scale} gives
$|x_j^p(t)|\le R_x+TV_\eps$. Choose a fixed polylogarithmic
$Q_\eps\ge1$ dominating $|q|$ on that actual phase region, independently
of $R$. Use the globally bounded test
$\clip(q^{>R}/Q_\eps,-1,1)$ on every signed branch. On actual endpoints
restoring $Q_\eps$ gives the raw tail, and on artificial branches its
absolute value is bounded by $|q^{>R}|$. Restore degree two only after
H3's unweighted remainder and tail estimates. Choose their arbitrary-power
accuracy after the degree, requested $M$, and the covariance factor $\mu$.
All restored errors are $O(\eps^M)$ uniformly in $R$.

Every tested endpoint is on a required record-bearing line continued to zero.
Its initial input is supported in the compact position support of $f_0$,
either through $F_{1,0}^\eps$ or through an initial cumulant dominated by
products of $f_0$. Retain this anchor before spatial envelope domination.
H4 gives, on every nonzero retained branch with birth energy $E$,
\[
 |v_j(t)|\le\sqrt E,\qquad |x_j(t)|\le R_x+T\sqrt E.
\]
The continued line is position-continuous. Thus for $R\ge\max(1,2R_x)$
a nonzero tail factor forces
$E\ge R^2/[4(T+1)^2]$.
If $q$ has degree $d$, two raw endpoint factors are bounded by
\[
 C_{q,T}(1+E)^d\ind_{\{E\ge c_TR^2\}}.
\]
Only after extracting this implication discard the compact-anchor indicator.
One birth factor is integrated in the common translation and the others
bounded by spatial suprema, leaving $C^m e^{-bE}$ as in H4. One fixed
Gaussian fraction extracts $e^{-c_qR^2}$, a second absorbs the fixed
polynomial, and the rest supports the same factorial and time sum. The
gamma-ratio estimate in Lemma~\ref{lem:weighted} shows that fixed degree
changes only the size polynomial. The reserved fractions and the admissible
size base are fixed once; $L$ is not chosen after $R$ or $q$.
At connected order $h=1,2$ the bound is
$C_q\mu^{1-h}e^{-c_qR^2}$, uniformly in record times. Later independent
births need not have compact spatial support.

For a field test $u$ the exact second-cumulant identity is
\[
 \operatorname{Var}\zeta_t^{\eps,p}(u)
       =\int F_{1,t}^{\eps,p}u^2
                          +\mu\int c_{2,t}^{\eps,p}(u\otimes u).
\]
The retained bounds at $u=q^{>R}$ and the restored errors prove
\eqref{eq:endpoint-tail}. The same proof without a tail gives bounded
second moments. Cauchy--Schwarz then proves two-time covariance cutoff
approximation, uniformly in its two record times.
\end{proof}

For $q\in\cQ$, the true law inherits a marginal probability cutoff gate.
Indeed $|q^{>R}|\le|q|$, so \eqref{eq:global-bound} and the Holder estimate
\eqref{eq:mean-transfer} compare the exact centerings uniformly in $R,t$.
Off H6's full-time exceptional event the raw paths agree. Chebyshev applied
to \eqref{eq:endpoint-tail}, then this comparison, gives for every $\delta>0$
\[
 \lim_{R\to\infty}\limsup_{\eps\to0}\sup_{t\le T}
           \Prob\{|\zeta_t^\eps(q^{>R})|>\delta\}=0.
\]
It is a marginal statement, not an assertion of a true covariance bound.

Uniform marginal variance alone does not imply a uniform path cutoff. For
example, a randomly placed interval of length $1/m$ on the unit time circle
has an indicator with variance at most $1/m$ at each fixed time, whereas its
centered time supremum is $1-1/m$. Theorem~\ref{thm:main} avoids that issue
by \eqref{eq:exact-path} and the monotone speed-sum LLN.
Likewise taking $R$ beyond the actual polylog endpoint range at each fixed
$\eps$ would produce a vacuous tail bound; the iterated order stated in the
lemma is the relevant one.

For $q\in\cQ$ the lemma also defines the Gaussian $L^2$ endpoint limit
of $G_t(q\chi_R)$. Apply its bound to differences of compact cutoffs,
pass to their fixed finite-dimensional Gaussian limits, and obtain an
$L^2$-Cauchy family. Joint probability comparison with the moment limit
already proved identifies it as $Z_0(r_q(t))$. This supplies the indicated
finite-dimensional covariance interpretation, not an all-polynomial
strong-dual functional theorem.

\subsection{The unresolved covariance transfer}
On the bad event the mechanical bound allows a quadratic fluctuation
difference of order $\eps^{-1}N/\mu$. H1 and H6 give only
\begin{equation}\label{eq:failed-L2-transfer}
 \eps^{-2}\E[(N/\mu)^2\ind_{E_\eps}]
       \le C_p\eps^{-2+(21/20)(1-2/p)}.
\end{equation}
Its exponent approaches $-19/20$ as $p\to\infty$ and equals $-1$ at
$p=42$. Thus this method does not establish true prelimit covariance
convergence. A growing upper bound supplies no lower bound on the actual
covariance and proves no physical divergence.

A scalar example displays the logical gap. Take a bounded centered $W$
and an independent event $A_\eps$ of probability $\eps^{21/20}$, and set
\[
 X_\eps=W+\eps^{-1}(\ind_{A_\eps}-\eps^{21/20}).
\]
Then $X_\eps-W\to0$ in probability, its uncentered mean shift
$\eps^{1/20}$ vanishes, but
\[
 \operatorname{Var}X_\eps-\operatorname{Var}W
        =\eps^{-19/20}(1-\eps^{21/20})\to\infty.
\]
This is not a hard-sphere counterexample. Additional weighted true-flow
control could prove covariance convergence; the present results leave it
unresolved.

\subsection{Other boundaries}
The theorem identifies global quadratic moments and their first means.
It does not establish a phase-space distribution limit of
$\eps^{-1}(\E\pi_t^\eps-f_t)$, a full unit-activity mean corrector,
or a local stress-corrector equation. No compact cubic activity result
is used or reproduced. The invariant projection and Gaussian mapping are
standard; the additional work is the clipped two-root weighted record
bound, its uniform weak-law use, the exact conserved closure, and the
intrinsic corollary, all relative to the pinned package.

The compact initial position anchor is substantive for
Lemma~\ref{lem:endpoint}. With large initial position and small velocity,
its implication $E\ge cR^2$ fails. A summable spatial Bol norm alone
does not force a second spatial moment. The compact support assumption
also makes the initial polynomial-to-compact substitution exact.
Removal of this assumption requires separate weighted initial information.
